\documentclass[11pt,a4paper]{article}
\usepackage[utf8]{inputenc}
\usepackage[spanish]{babel}
\usepackage{amsmath}
\usepackage{amsfonts}
\usepackage{amssymb}
\usepackage{graphicx}
\usepackage{booktabs}
\usepackage{siunitx}
\usepackage{caption}
\usepackage{subcaption}
<<<<<<< HEAD
\usepackage{placeins}
=======
>>>>>>> c4b7d17 (supercipotes)
\usepackage{geometry}
\geometry{margin=2.5cm}

\sisetup{
  separate-uncertainty = true,
  per-mode = symbol
}

% Datos de la práctica
\def\numeroPractica{10}
\def\tituloPractica{Característica resistencia-temperatura de un termistor}

\title{Práctica \numeroPractica: \tituloPractica}
<<<<<<< HEAD
\author{Andrés Vinuesa Espinosa y José María Martínez Herrada}
=======
\author{Andrés}
>>>>>>> c4b7d17 (supercipotes)
\date{\today}

\begin{document}

\maketitle

\begin{abstract}
En esta práctica se estudia la dependencia de la resistencia eléctrica de un termistor con la temperatura absoluta. A partir de medidas experimentales, se verifica el modelo exponencial $R = A e^{B/T}$. Los parámetros característicos del termistor se obtienen mediante dos aproximaciones: un cálculo analítico a partir de dos puntos de referencia y un ajuste no lineal por regresión de distancia ortogonal (ODR) considerando incertidumbres en ambas variables. El ajuste ODR proporciona los valores $A = (1.59 \pm 0.08) \times 10^{-5} \text{ k}\Omega$ y $B = 3973 \pm 14 \text{ K}$, con un excelente coeficiente de bondad de ajuste, confirmando la validez del modelo teórico asumido.
\end{abstract}

\section{Introducción y Objetivos}
% [HUMAN AUTHOR: INTRODUCCIÓN AQUÍ]
El objetivo de esta práctica es determinar la relación experimental entre la resistencia eléctrica de un termistor y su temperatura absoluta, ajustando los datos a un modelo exponencial de la forma $R = A e^{B/T}$.

\section{Resultados}

\subsection{Cálculo analítico a partir de dos puntos}
Se han seleccionado dos puntos de referencia para realizar un cálculo preliminar de los parámetros $A$ y $B$:
\begin{itemize}
<<<<<<< HEAD
    \item $T_1 = 353.15 \pm 0.10\,\mathrm{K}$, $R_1 = 1.25 \pm 0.22\,\mathrm{k\Omega}$
    \item $T_2 = 273.25 \pm 0.10\,\mathrm{K}$, $R_2 = 31.50 \pm 0.58\,\mathrm{k\Omega}$
=======
    \item $T_1 = \qty{353.15 \pm 0.10}{K}$, $R_1 = \qty{1.25 \pm 0.22}{k\Omega}$
    \item $T_2 = \qty{273.25 \pm 0.10}{K}$, $R_2 = \qty{31.50 \pm 0.58}{k\Omega}$
>>>>>>> c4b7d17 (supercipotes)
\end{itemize}

A partir de la solución analítica del sistema de ecuaciones:
\begin{equation}
    B = \ln\left(\frac{R_1}{R_2}\right) \frac{T_1 T_2}{T_2 - T_1} \quad ; \quad A = R_1 e^{-B/T_1}
\end{equation}
Se han obtenido los siguientes valores:
\begin{itemize}
<<<<<<< HEAD
    \item $A = 2.0 \pm 1.6\times 10^{-5}\,\mathrm{k\Omega}$
    \item $B = 3890 \pm 210\,\mathrm{K}$
=======
    \item $A = \qty{2.0 \pm 1.6e-5}{k\Omega}$
    \item $B = \qty{3890 \pm 210}{K}$
>>>>>>> c4b7d17 (supercipotes)
\end{itemize}
Este método proporciona una estimación inicial de los parámetros, pese a que es ampliamente suscepctible a errores, debido a que
solo se toman dos medidas, este método proporciona una buena estimación de parámetros para así obtener un buen punto de partida para el algoritmo
en la siguiente sección, el cual es más robusto al considerar todas las medidas experimentales.

\subsection{Representación gráfica y Ajuste ODR}
Se ha representado el logaritmo de la resistencia en función de la inversa de la temperatura absoluta ($1/T$), lo cual debería mostrar una tendencia lineal según la ecuación $\ln(R) = \ln(A) + B/T$.

\begin{figure}[ht!]
    \centering
    \includegraphics[width=0.8\textwidth]{logvst.png}
    \caption{Representación de $\ln(R)$ frente a $1/T$ ($K^{-1}$).}
    \label{fig:logvst}
\end{figure}

<<<<<<< HEAD
Los datos experimentales utilizados en el ajuste se recogen en la Tabla~\ref{tab:datos_experimentales} del anexo.

Posteriormente, se ha realizado un ajuste no lineal mediante el método de Ortogonal Distance Regression (ODR) sobre la relación $R = A e^{B/T}$, considerando errores en ambas variables. Los parámetros obtenidos del ajuste son:
\begin{itemize}
    \item $A = 1.59 \pm 0.08\times 10^{-5}\,\mathrm{k\Omega}$
    \item $B = 3973 \pm 14\,\mathrm{K}$
=======
Posteriormente, se ha realizado un ajuste no lineal mediante el método de Ortogonal Distance Regression (ODR) sobre la relación $R = A e^{B/T}$, considerando errores en ambas variables. Los parámetros obtenidos del ajuste son:
\begin{itemize}
    \item $A = \qty{1.59 \pm 0.08e-5}{k\Omega}$
    \item $B = \qty{3973 \pm 14}{K}$
>>>>>>> c4b7d17 (supercipotes)
    \item Coeficiente de determinación $R^2 = 0.999876$
\end{itemize}


Dado que $R^2 \approx 1$, se puede concluir que existe una correlación entre ambas variables, sin embargo, no podemos concluir que esta sea la correcta,
ya que es posible que haya factores externos que no modifiquen la correlación, pero sí modifiquen la relación entre ambas variables, por lo que es importante analizar los resultados con cautela.
Por ejemplo, una de las posibles limitaciones del modelo teórico es que no considera temperaturas extremas, muy bajas o muy altas, que podrían alterar la estructura misma del material,
lo cual podría modificar esta relación entre resistencia y temperatura, por ende es de vital importancia considerar que los resultados obtenidos en esta práctica
solamente son válidos para el rango de temperaturas estudiado, y no se pueden extrapolar a temperaturas extremas sin realizar un análisis más detallado.
\subsection{Medición de temperatura}
<<<<<<< HEAD
Al comparar los parámetros característicos obtenidos en los apartados anteriores, observamos que los valores son compatibles entre sí, ya que se solapan dentro de los márgenes de error. Sin embargo, el ajuste ODR proporciona una precisión mucho mejor. Para el parámetro $A$, el cálculo analítico da un valor de $A = (2,0 \pm 1,6) \times 10^{-5} \text{ k}\Omega$, con un error relativo altísimo, mientras que el ajuste nos da un resultado de $A = (1,59 \pm 0,08) \times 10^{-5} \text{ k}\Omega$. Lo mismo ocurre con el parámetro $B$, que pasa de $B = (3890 \pm 210) \text{ K}$ en el cálculo de dos puntos a $B = (3973 \pm 14) \text{ K}$ en el ajuste. Esto nos confirma la mayor precisión que tiene el ajuste respecto al método analítico.

Para mirar si son buenos nuestros resultados, sumergiremos el termistor en un baño de agua a temperatura ambiente y mediremos su resistencia. Usaremos la siguiente ecuación para determinar la temperatura del agua a partir de dicha resistencia:

\begin{equation}
T = \frac{B}{\ln(R/A)}
\end{equation}

Aplicando esta ecuación con las dos parejas de constantes calculadas, obtenemos dos estimaciones para la temperatura del agua:
\begin{itemize}
    \item Con los parámetros del cálculo de dos puntos se obtiene una temperatura de $(21,4 \pm 2,3) \text{ }^{\circ}\text{C}$.
    \item Con los parámetros del ajuste ODR se obtiene una temperatura de $(22,35 \pm 1,50) \text{ }^{\circ}\text{C}$.
\end{itemize}

Ambos valores son razonables y coherentes con la temperatura ambiente habitual del laboratorio. No obstante, el valor calculado a partir de los parámetros del ajuste presenta una incertidumbre bastante menor, demostrando que el cálculo mediante el ajuste resulta en un termómetro más preciso y fiable.

\section{Conclusiones}
En esta práctica se ha estudiado la exponencial de la resistencia eléctrica de un termistor con respecto a la temperatura absoluta, siguiendo el modelo teórico $R=Ae^{B/T}$. Se han contrastado dos metodologías: un cálculo analítico preliminar basado en dos puntos de referencia y un ajuste no lineal mediante regresión de distancia ortogonal (ODR).

Los resultados obtenidos demuestran que es mejor el método de ajuste en términos de precisión. Mientras que el cálculo analítico a partir de dos puntos arrojó márgenes de error muy elevados debido a la escasez de datos, el algoritmo ODR permitió determinar con gran precisión las constantes características del termistor, obteniendo los valores $A=(1,59\pm0,08)\times10^{-5}\text{ k}\Omega$ y $B=(3973\pm14)\text{ K}$. El elevado coeficiente de Pearson obtenido ($R^2=0.999876$) confirma la excelente bondad del ajuste y valida el modelo exponencial asumido para el rango de trabajo. En cambio el método analítico nos da $A=(2.0\pm1.6)\times10^{-5}\text{ k}\Omega$ y $B=(3890\pm210)\text{ K}$ que tienen unas incertidumbres altísimas, lo que nos dice que el ajuste nos da unos resultados más fiabels.

Por otro lado, la aplicación práctica de estas constantes en la ecuación termométrica para determinar la temperatura del agua del grifo ratificó la fiabilidad del ajuste. Aunque ambos métodos arrojaron resultados compatibles con la temperatura ambiente del laboratorio, los parámetros del ajuste proporcionaron una medida notablemente más precisa y con menor incertidumbre experimental ($(22.4\pm1.5)\text{ }^{\circ}\text{C}$). Usando las constantes obtenidas por el otro método nos da una temperatura de $21.4\pm2.3$ que es compatible con la temperatura ambiente, pero con una incertidumbre mucho mayor.


=======
Para contrastar nuestros hallazgos, sumergiremos el termistor en un baño de agua a temperatura ambiente, y mediremos su resistencia, para así comparar el valor medido con el valor predicho por el modelo, y así verificar la validez del mismo.
usaremos la siguiente ecuación para calcular la temperatura del agua a partir de la resistencia medida:
\begin{equation}
    T = \frac{B}{\ln(R/A)}
\end{equation}
Una vez realizado el cálculo, se ha obtenido una temperatura de \qty{21.4 \pm 2.3}{\degreeCelsius}, lo cual corresponde a una temperatura de \qty{22.35 \pm 1.5}{\degreeCelsius}, lo cual es razonable para la temperatura ambiente del laboratorio.

\section{Conclusiones}
% [HUMAN AUTHOR: ESCRIBE LAS CONCLUSIONES AQUÍ]
% Resume los aspectos más relevantes y reflexiona sobre los resultados obtenidos.
Esto es una geminada chema, pon tu algo
>>>>>>> c4b7d17 (supercipotes)
\appendix
\section{Anexos}

\subsection{Tablas de datos}
<<<<<<< HEAD

\begin{table}[ht!]
    \centering
    \begin{tabular}{r r r r r}
    \toprule
    {Temp. ($^{\circ}\mathrm{C}$)} & {Temp. (K)} & {$u_T$ (K)} & {Resistencia ($k\Omega$)} & {$u_R$ ($k\Omega$)} \\
    \midrule
    5.0 & 278.15 & 0.10 & 25.25 & 0.50 \\
    20.0 & 293.15 & 0.10 & 12.26 & 0.35 \\
    25.0 & 298.15 & 0.10 & 9.81 & 0.32 \\
    30.0 & 303.15 & 0.10 & 7.80 & 0.29 \\
    35.0 & 308.15 & 0.10 & 6.32 & 0.28 \\
    40.0 & 313.15 & 0.10 & 5.13 & 0.26 \\
    45.0 & 318.15 & 0.10 & 4.20 & 0.25 \\
    50.0 & 323.15 & 0.10 & 3.46 & 0.24 \\
=======
En la Tabla~\ref{tab:datos_experimentales} GEMINADAAA REVISAR se muestran los valores medidos de temperatura y resistencia con sus respectivas incertidumbres.

\begin{table}[ht!]
    \centering
    \begin{tabular}{S[table-format=2.1] S[table-format=3.2] S[table-format=1.2] S[table-format=2.2] S[table-format=1.3]}
    \toprule
    {Temp. (\unit{\degreeCelsius})} & {Temp. (\unit{K})} & {$u_T$ (\unit{K})} & {Resistencia (\unit{k\Omega})} & {$u_R$ (\unit{k\Omega})} \\
    \midrule
    5.0 & 278.15 & 0.10 & 25.25 & 0.503 \\
    20.0 & 293.15 & 0.10 & 12.26 & 0.347 \\
    25.0 & 298.15 & 0.10 & 9.81 & 0.318 \\
    30.0 & 303.15 & 0.10 & 7.80 & 0.294 \\
    35.0 & 308.15 & 0.10 & 6.32 & 0.276 \\
    40.0 & 313.15 & 0.10 & 5.13 & 0.262 \\
    45.0 & 318.15 & 0.10 & 4.20 & 0.250 \\
    50.0 & 323.15 & 0.10 & 3.46 & 0.242 \\
>>>>>>> c4b7d17 (supercipotes)
    \bottomrule
    \end{tabular}
    \caption{Datos experimentales de temperatura y resistencia medidos.}
    \label{tab:datos_experimentales}
\end{table}
<<<<<<< HEAD
\FloatBarrier

\subsection{Análisis de Incertidumbres}

En este apartado se evalúan las incertidumbres asociadas a las magnitudes medidas directamente en el laboratorio, así como la propagación sistemática de errores para aquellas variables obtenidas de forma indirecta a través de relaciones matemáticas.

\subsubsection{Medidas Directas}

\begin{itemize}
    \item \textbf{Temperatura ($T$):} La temperatura del baño térmico se registró mediante un termómetro digital de precisión. La incertidumbre asociada se considera de tipo B (instrumental), determinada por la resolución del propio dispositivo, estableciéndose un valor constante de $u_T = 0,10\text{ K}$ para todo el rango de trabajo.
    
    \item \textbf{Resistencia eléctrica ($R$):} Las medidas de resistencia se realizaron utilizando un polímetro digital \textit{PeakTech 2005} \cite{manual_peaktech}. La incertidumbre instrumental de tipo B se determinó a partir de las especificaciones técnicas del fabricante proporcionadas en el manual del usuario. El error se define mediante una combinación lineal de un porcentaje de la lectura y un número determinado de dígitos del contador. Debido a este criterio, la incertidumbre absoluta $u_R$ varía de forma dinámica en cada punto experimental según la escala y la magnitud de la resistencia medida, tal y como queda reflejado en la Tabla~\ref{tab:datos_experimentales}.
\end{itemize}

\subsubsection{Medidas Indirectas}

Para las magnitudes calculadas de manera indirecta, la incertidumbre combinada se obtiene mediante la ley de propagación cuadrática de errores, asumiendo que las covarianzas entre las variables independientes son nulas.

\paragraph{1. Parámetros $A$ y $B$ mediante el cálculo de dos puntos:}
A partir de los puntos de referencia $(T_1, R_1)$ y $(T_2, R_2)$, el parámetro $B$ se calcula mediante la relación indirecta:
\begin{equation}
B = \ln\left(\frac{R_1}{R_2}\right) \frac{T_1 T_2}{T_2 - T_1}
\end{equation}

La propagación de su incertidumbre viene dada por:
\begin{equation}
u_B = \sqrt{\left(\frac{\partial B}{\partial T_1} u_{T_1}\right)^2 + \left(\frac{\partial B}{\partial T_2} u_{T_2}\right)^2 + \left(\frac{\partial B}{\partial R_1} u_{R_1}\right)^2 + \left(\frac{\partial B}{\partial R_2} u_{R_2}\right)^2}
\end{equation}

Donde las derivadas parciales analíticas correspondientes son:
\begin{align}
\frac{\partial B}{\partial R_1} &= \frac{T_1 T_2}{R_1(T_2 - T_1)} & \frac{\partial B}{\partial R_2} &= -\frac{T_1 T_2}{R_2(T_2 - T_1)} \\
\frac{\partial B}{\partial T_1} &= \ln\left(\frac{R_1}{R_2}\right) \frac{T_2^2}{(T_2 - T_1)^2} & \frac{\partial B}{\partial T_2} &= -\ln\left(\frac{R_1}{R_2}\right) \frac{T_1^2}{(T_2 - T_1)^2}
\end{align}

Para el parámetro $A$, expresado como $A = R_1 e^{-B/T_1}$, su incertidumbre se propaga de forma análoga a partir del desglose algebraico de sus variables base, resultando en las siguientes derivadas parciales analíticas directas:
\begin{align}
\frac{\partial A}{\partial R_1} &= A \left( -\frac{T_1}{T_2 - T_1} \frac{1}{R_1} \right) & \frac{\partial A}{\partial R_2} &= A \left( \frac{T_2}{T_2 - T_1} \frac{1}{R_2} \right) \\
\frac{\partial A}{\partial T_1} &= -A \frac{T_2}{(T_2 - T_1)^2} \ln\left(\frac{R_1}{R_2}\right) & \frac{\partial A}{\partial T_2} &= A \frac{T_1}{(T_2 - T_1)^2} \ln\left(\frac{R_1}{R_2}\right)
\end{align}

\paragraph{2. Determinación de la temperatura del agua del grifo:}
La medición indirecta de la temperatura $T$ a partir de la ecuación termométrica depende de las constantes calibradas ($A, B$) y de la resistencia medida ($R$):
\begin{equation}
T = \frac{B}{\ln(R/A)}
\end{equation}

La expresión para la propagación de la incertidumbre en este caso es:
\begin{equation}
u_T = \sqrt{\left(\frac{\partial T}{\partial B} u_B\right)^2 + \left(\frac{\partial T}{\partial R} u_R\right)^2 + \left(\frac{\partial T}{\partial A} u_A\right)^2}
\end{equation}

Cuyas derivadas parciales analíticas se desarrollan de la siguiente forma simplificada:
\begin{align}
\frac{\partial T}{\partial B} &= \frac{1}{\ln(R/A)} = \frac{T}{B} \\
\frac{\partial T}{\partial R} &= -\frac{B}{R \cdot [\ln(R/A)]^2} = -\frac{T^2}{B \cdot R} \\
\frac{\partial T}{\partial A} &= \frac{B}{A \cdot [\ln(R/A)]^2} = \frac{T^2}{B \cdot A}
\end{align}

Sustituyendo los valores numéricos experimentales y las incertidumbres de los parámetros en estas expresiones, se obtienen las desviaciones finales reportadas en la sección de resultados para cada método de calibración.


\begin{thebibliography}{9}

\bibitem{guion_practicas}
Departamento de Física Aplicada, \emph{Guion de Prácticas de Termodinámica: Práctica 10 - Calibrado de termistores}. Universidad de Granada, Curso 2025-2026.

\bibitem{manual_peaktech}
PeakTech, \emph{Manual de instrucciones Multímetro Digital PeakTech 2005}. Disponible en línea: \texttt{https://www.denver.es/wp-content/uploads/2018/01/PeakTech2005100817ES.pdf}

\end{thebibliography}

=======

\subsection{Análisis de Incertidumbres}
hacer esto me ha dado una pereza q flipas xd menciona el manual este del polimetro
https://www.denver.es/wp-content/uploads/2018/01/PeakTech_2005_100817_ES.pdf
y saca las cosas de ahí, como la incertidumbre de la resistencia medida, que es el 0.08 del valor medido, y la incertidumbre de la temperatura, que es de 0.1 K, y luego menciona que estas incertidumbres se han propagado a través de las fórmulas para obtener las incertidumbres finales de los parámetros A y B, así como de la temperatura calculada del baño de agua.
>>>>>>> c4b7d17 (supercipotes)

\end{document}
